\section{How each figure was made}
\label{sec:provenance}

This section describes how every figure and table in the report was produced,
so that a reader can judge what each one can support. The descriptions are of
what was computed and from what. They name no software, because the report
describes the experiment rather than the code that ran it.

One principle governs all of them. Every figure is regenerated from the result
files the experiments wrote, and no number is transcribed by hand at any point.
A figure therefore cannot drift away from the run it describes: when a policy is
retrained, the figures are redrawn and the numbers follow. The tables are
produced the same way and inserted whole, so a value in a table and the same
value in the text cannot disagree.

\paragraph{The recordings, Figure~\ref{fig:contact}.}
One frame is taken from the temporal midpoint of every recording, using the
overhead camera, and the frames are laid out in recording order. The midpoint
rather than the opening frame, because every recording begins with the same
flattened garment and a sheet of opening frames would be sixty-five copies of
one picture. The outline marks the recordings the trainer withheld, read from
the trained model's own configuration rather than assumed, so the figure cannot
disagree with the split that was actually used.

\paragraph{The action error, Table~\ref{tab:policies} and
Figures~\ref{fig:gap} and~\ref{fig:horizon}.}
Every twentieth frame of a recording is given to the policy as an observation.
The policy plans a chunk of future commands, and that chunk is compared step by
step against what the demonstrator actually commanded next. Squaring and
averaging those differences over all sampled frames, separately at each position
within the chunk, gives the curve in Figure~\ref{fig:horizon}; averaging that
curve over its whole length and taking the square root gives one number per
policy. The shared-horizon column is the same curve averaged over only its first
thirty-two positions, which is why the two columns can differ for the same
policy and why both are reported. Policies whose output is sampled have their
sampler fixed to a recorded seed, so that rerunning the scoring reproduces the
number rather than redrawing it. The fine-tuned family reads a written
instruction as well as its cameras, and it is scored with the instruction it
was trained on; an earlier scoring that supplied none put its held-out error
about a tenth higher.

\paragraph{The input contributions, Figure~\ref{fig:streams}.}
For each sampled frame the policy plans once with all inputs present. Each input
group is then replaced in turn by an uninformative substitute --- the average of
that input across the dataset, so the substitution carries no information rather
than introducing an artificial edge --- and the policy plans again. The size of
the change between the two plans is that group's contribution on that frame.
Contributions are summed over frames and normalised to sum to one within each
policy, which is what restricts the figure to statements about ordering.

\textbf{A correction to this figure.} An earlier version drew the fine-tuned
family's bar from a run in which its trained adapter had not been loaded, so it
described an untrained model; that bar put proprioception at a twentieth and
touch at two fifths. The bar shown here reruns the same checkpoint on the same
recordings with the adapter loaded and the task instruction supplied. The other
two bars were never affected, because those families store complete weights.

\paragraph{The rim measure, Figure~\ref{fig:rim}.}
Gradient maps are computed for the tactile cameras, giving one map per camera
per frame. A band is defined around the edge of each map in \emph{normalised}
coordinates --- a fraction of the image, not a count of cells --- because the
policy families produce maps at different resolutions and a fixed number of
cells would mean a different fraction of the image for each. The mass falling
inside the band is divided by the mass a completely flat map of the same
resolution would place there, which is what makes a value comparable across
resolutions and gives the dashed reference line its meaning. Three band widths
are reported so that the conclusion cannot depend on one arbitrary choice.

\paragraph{The hold-out control, Table~\ref{tab:split} and
Figure~\ref{fig:split}.}
Each family is trained a second time, changing one thing: instead of withholding
the last seven recordings, seven are withheld at random from across the session.
Architecture, optimisation budget, random seed and scoring are all held at the
values the first training used, so the two runs of a family differ in which
recordings were withheld and in nothing else. Each trained model is then scored
on the recordings it saw and on the recordings it did not, and the ratio of the
two is the gap. A ratio is used rather than a difference because it is
dimensionless and computed within one policy, which is what allows the two
families to appear on one axis despite planning different distances ahead. The
per-recording figures quoted alongside come from scoring each withheld
recording on its own instead of pooling the seven.

\paragraph{The framewise comparison, Figures~\ref{fig:framewise}
and~\ref{fig:gradcampeak}, and the videos.}
All four are computed on one recording the policies never saw, sampling every
fifth frame of it. At each sampled frame every policy plans, each input is
suppressed in turn exactly as described above, and the resulting contributions
are kept per frame instead of being pooled. Figure~\ref{fig:framewise} sums the
four fingertip cameras at each instant and plots that against time from the
start of the recording. Figure~\ref{fig:gradcampeak} is one instant of the
gradient picture, chosen as the frame where the contribution of touch reaches
its maximum, and it is rendered directly rather than recovered from a compressed
video so that the panels are at full quality.

Two rules are enforced when the panels are assembled rather than left to care.
The comparison covers only the frames that every policy scored, so an
incomplete run shortens it instead of putting different policies at different
moments. And a cropped policy's tactile panels are produced by applying that
policy's own recorded crop to the frame, so the panel shows the image the policy
received; the crop is read from the trained model's configuration rather than
restated, which is what stops the picture and the policy disagreeing.

The two moving versions accompany the report as video files rather than being
embedded, because the document format cannot carry them: one shows the gradient
pictures through the recording, the other the contribution of every input
through the same recording at the same instants. A policy with no spatial
feature map is left out of the gradient video by name rather than drawn as an
empty panel, since an empty panel beside full ones reads as a policy attending
to nothing.
